% TOP_PROBLEM: {"id": 8400005, "title": "O5b — Randomized polynomial-time integer factorization", "category_id": 3, "category": {"id": 3, "name": "graph_theory", "display_name": "Graph Theory", "description": "Problems involving graphs, networks, and their properties.", "slug": "graph-theory", "order_index": 3, "created_at": "2026-07-31T15:26:25.671Z"}, "status": "open", "problem_number": "AMR-083-0005", "set_id": 13}
% FIRST_POSTED: 2026-10-01
% ENTRY_KIND: statement_only
% UnsolvedMath ID 8400005; source catalogue code AMR-083-0005
% !TeX program = lualatex
% Definition and sources only; this file is not an LLM solution attempt.
\documentclass[11pt,a4paper]{article}
\usepackage[margin=25mm]{geometry}
\usepackage{fontspec}
\setmainfont{lmroman10-regular.otf}[BoldFont=lmroman10-bold.otf,
 ItalicFont=lmroman10-italic.otf,BoldItalicFont=lmroman10-bolditalic.otf]
\usepackage{amsmath,amssymb,mathtools}
\usepackage{unicode-math}
\setmathfont{latinmodern-math.otf}
\AtBeginDocument{\let\setminus\smallsetminus}
\usepackage{xurl,hyperref}
\hypersetup{colorlinks=true,urlcolor=blue}
\setcounter{secnumdepth}{0}
\setlength{\parindent}{0pt}
\setlength{\parskip}{5pt}
\setlength{\emergencystretch}{3em}
\begin{document}
\section*{O5b — Randomized polynomial-time integer factorization}\label{8400005}
{\small UnsolvedMath ID 8400005 · AMR-083-0005 · Graph Theory · AMR Open Problem Lists}
\subsection{Definitions and mathematical statement}
For an integer $N\ge2$, its complete prime factorization is the list of distinct primes and positive exponents
\[
N=p_1^{e_1}\cdots p_t^{e_t},\qquad p_1<\cdots<p_t.
\]
The input and every integer in the output are written in binary. The problem includes arbitrary composite inputs and prime powers, not only products of two primes.

Does a classical randomized algorithm compute this factorization in polynomial time in $\ell(N)=\lfloor\log_2N\rfloor+1$? An explicit equivalent search formulation asks for constants $C,c>0$ and a probabilistic algorithm which, on every $N$, uses at most $C(1+\ell(N))^c$ bit operations and returns the correct complete list with probability at least $2/3$, the probability being over its internal random bits. The guarantee is required separately for each input, rather than on average over a distribution of integers.

One can instead require a Las Vegas algorithm: every returned list is correct and the expected running time is polynomial in $\ell(N)$. These formulations agree here because a proposed list can be checked by multiplying the prime powers and testing the listed factors for primality; failed or incorrect trials can be rejected and repeated. Conversely, truncating an expected-polynomial computation gives constant success probability. The target is classical bit complexity; access to quantum operations is not part of the computational model.
\subsection{Short English statement}
Does complete integer factorization admit a classical randomized algorithm with polynomial running time?
\subsection{Sources}
\begin{itemize}
\item[S1] ULAM.AI and the UnsolvedMath Contributors, supplied problems.json, record 8400005 (AMR-083-0005), AMR Open Problem Lists. Catalogue identity and source transcription; CC BY 4.0.
\url{https://huggingface.co/datasets/ulamai/UnsolvedMath}.
\item[S2] Adleman \& McCurley - Open Problems in Number Theory Complexity II (1994), O5b, p7 / LNCS p297. Original source indexed by AMR-083-0005.
\url{https://mccurley.org/papers/open.ps.gz}.
\end{itemize}
\end{document}
